\section{Localisation $\Reop(s) = 1/2$ : quatre mécanismes équivalents}
\label{sec:localisation_critique}

L'objet central de cette section est l'explication PT-canonique de la
ligne critique. Quatre mécanismes apparemment distincts apparaissent à des
étapes différentes du programme. Ils sont en fait quatre faces d'une
unique structure~: la mesure invariante de la dynamique de persistance.

\subsection{Cellule de Planck symplectique : $u_\star\,p_\star = 2\pi$}
\label{sec:cellule_planck}

L'opérateur de dilatation $\HPTBK$ engendre les transformations d'échelle
sur la coordonnée de Mellin. Sa mesure de Liouville canonique est
$du\,dp_u / (2\pi)$, conséquence directe du commutateur canonique
$[u, p_u] = i$.

\begin{proposition}[Cellule de Planck PT]
\label{prop:cellule_planck}
La cellule unité de l'espace de phases canonique est
\begin{equation}
\boxed{\;u_\star\, p_\star \;=\; 2\pi\;}
\label{eq:cellule_planck}
\end{equation}
avec le choix symétrique
$u_\star = p_\star = \sqrt{2\pi} \approx 2{,}5066$ (point fixe de
l'involution $(u, p_u) \leftrightarrow (p_u, u)$). À énergie $\gamma$
fixée, l'hyperbole $\HPTBK = \gamma$ ne pénètre dans l'espace admissible
que sur le segment
\[
u \,\in\, [u_\star,\, \gamma/p_\star] \;=\; [\sqrt{2\pi},\, \gamma/\sqrt{2\pi}],
\]
et le cut-off dynamique canonique vaut
\begin{equation}
\umax(\gamma) \;=\; \dfrac{\gamma}{\sqrt{2\pi}}.
\label{eq:umax}
\end{equation}
\end{proposition}

\begin{proposition}[Densité semi-classique microcanonique]
\label{prop:densite_NPT}
Avec le cut-off~\eqref{eq:umax} et la double barrière imposée par
l'involution $(u, p_u) \leftrightarrow (p_u, u)$, la quantification
microcanonique de $\HPTBK$ en ordre symétrique de Weyl produit la densité
intégrale
\begin{equation}
\NPT(\gamma) \;=\; \dfrac{\gamma}{2\pi}\,\log\dfrac{\gamma}{2\pi\,e} \,+\, 1.
\label{eq:NPT}
\end{equation}
\end{proposition}
\label{sec:densite_semi_classique}

\begin{proof}[Esquisse]
L'aire microcanonique entre la courbe $\HPTBK = \gamma$ et la barrière
$u = u_\star$ dans le quadrant $u, p_u > 0$ vaut
$A(\gamma) = \int_{u_\star}^{\umax(\gamma)} \gamma/u\,du
= \gamma\,\log\!\big(\umax(\gamma)/u_\star\big)
= \gamma\,\log\!\big(\gamma/2\pi\big)$.
La double barrière (le secteur miroir $u, p_u < 0$ obtenu par l'involution)
double cette aire. La quantification semi-classique de Weyl en ordre
symétrique donne $\NPT(\gamma) = A(\gamma)/\pi + \text{const.}$ La
constante est fixée à $+1$ par la double extension auto-adjointe de
$\HPTBK$ aux bords (voir \S\ref{sec:bord_antiperiodique}). Le calcul
algébrique direct donne~\eqref{eq:NPT}.
\end{proof}

\paragraph{Comparaison avec Riemann--von Mangoldt.}
La densité Riemann--von Mangoldt classique \cite{Edwards1974} est
\begin{equation}
N_{\mathrm{RvM}}(\gamma) \;=\; \dfrac{\gamma}{2\pi}\,\log\dfrac{\gamma}{2\pi\,e} \,+\, \dfrac{7}{8} \,+\, O(1/\gamma).
\label{eq:NRvM}
\end{equation}
La forme fonctionnelle de~\eqref{eq:NPT} et~\eqref{eq:NRvM} est
identique. L'écart constant
\begin{equation}
\NPT(\gamma) - N_{\mathrm{RvM}}(\gamma) \;=\; \dfrac{1}{8}
\label{eq:residu_un_huitieme}
\end{equation}
est vérifié à $10^{-15}$ sur quatre décades, $\gamma \in [14, 10^{3}]$.
Le $2\pi$ qui apparaît dans la densité $N(\gamma)$ est la mesure de
Liouville canonique sur $(u, p_u)$, conséquence du commutateur
$[u, p_u] = i$ ; le shift $\Reop(s) = 1/2$ est dérivé du Jacobien de la
mesure de Haar multiplicative (\S\ref{sec:haar_jacobien}). Le résidu
$1/8$ admet la sur-détermination cohomologique de
\S\ref{sec:surdetermination_un_huitieme}.

\subsection{Jacobien de Haar multiplicatif}
\label{sec:haar_jacobien}

Pour un mode stationnaire $f(\tau, u) = e^{i\omega\tau}\,\psi(u)$ de
l'opérateur d'Alembert sur la coordonnée de Mellin (la métrique Fisher
de $\Sigma_{\mathrm{pers}}$ étant lorentzienne pour
$\mu > \mu_c \approx 6{,}97$, \cite{PT_GeoSpectral}), le shift Mellin
sous la mesure de Haar multiplicative
$d\pi_{\log} = \sum_p \delta_{\log p}\,d\log p$ produit un facteur
$1/p = e^{-u/2}$ devant la transformée. Plus précisément, le pôle de la
transformée de Mellin se déplace de $s = \omega$ à
\begin{equation}
\boxed{\;s \;=\; \dfrac{1}{2} + i\omega.\;}
\label{eq:haar_shift}
\end{equation}
C'est le shift de la mesure de Haar multiplicative. La partie réelle
$1/2$ n'est pas postulée~: elle vient du Jacobien naturel de la mesure
invariante de Haar sur les premiers.

\subsection{Transformation unitaire $e^{su/2}$}
\label{sec:transformation_unitaire}

Le théorème d'auto-adjonction de Ruelle--PT (\cite{PT_Foundations})
établit que l'opérateur de Ruelle--PT discret $L_s$ est auto-adjoint pour
la mesure de Gibbs $\mu_s(p) = p^{-s}$ via la balance détaillée associée
à la symétrie antidiagonale $T_3 = T_3^{\top}$. En coordonnée de Mellin,
la mesure de Gibbs devient $\mu_s(p)\,dp = e^{-su}\,du$. La transformation
unitaire
\begin{equation}
U : f(u) \,\mapsto\, e^{-su/2}\,f(u)
\label{eq:U_unitaire}
\end{equation}
ramène l'opérateur $L_s$ sur $L^{2}(\Reals, e^{-su} du)$ à
$L^{2}(\Reals, du)$ en translatant $\HPTBK$ par $is/2$. Le shift apparaît
littéralement dans l'exposant~: c'est l'identification opératorielle
entre balance détaillée discrète ($\Ttwo$ du crible) et auto-adjonction
continue de $\HPTBK$.

\subsection{Condition aux bords antipériodique}
\label{sec:bord_antiperiodique}

Sur $L^{2}([\umin, \umax], du)$ avec $0 < \umin < \umax < \infty$, les
indices de défaut de $\HPTBK$ valent $(n_+, n_-) = (1, 1)$. Par von
Neumann, il existe une famille $U(1)$ d'extensions auto-adjointes
paramétrées par
\begin{equation}
\psi(\umax) \;=\; e^{i\theta}\,\psi(\umin), \qquad \theta \in [0, 2\pi).
\label{eq:BC_general}
\end{equation}

\begin{proposition}[Sélection PT-canonique : antipériodicité $\theta = \pi$]
\label{prop:antiperiodicite}
L'involution
\[
\iota : u \,\mapsto\, \umax + \umin - u
\]
est la projection continue de l'antidiagonale
$T_3 = \bigl(\begin{smallmatrix}0 & 1\\ 1 & 0\end{smallmatrix}\bigr)$ du
crible $\bmod\,3$. C'est l'unique involution continue qui préserve les
bords comme orbite et est une isométrie de $L^{2}$. Le terme de bord
$\iota\,\HPTBK\,\iota^{-1} + \HPTBK$ s'annule \textbf{uniquement} pour
$\theta = \pi$ (antipériodique), qui force la symétrie spectrale
$\{\gamma_n, -\gamma_n\}$.
\end{proposition}

\begin{proof}[Esquisse]
La famille~\eqref{eq:BC_general} est paramétrée par
$\theta \in [0, 2\pi)$. L'involution $\iota$ envoie une extension
auto-adjointe à paramètre $\theta$ sur celle à paramètre $-\theta$.
L'invariance sous $\iota$ impose donc $\theta \equiv -\theta \pmod{2\pi}$,
soit $\theta \in \{0, \pi\}$. Le cas $\theta = 0$ (périodicité) donne un
spectre symétrique mais ne préserve pas l'anti-symétrie de $\HPTBK$ aux
bords ; seul $\theta = \pi$ vérifie
$\iota\,\HPTBK\,\iota^{-1} + \HPTBK = 0$ et produit la symétrie spectrale
$\{\gamma_n, -\gamma_n\}$ requise. Vérifié numériquement à
$5 \cdot 10^{-12}$ (précision machine).
\end{proof}

L'antipériodicité continue est donc l'image continue de
l'antidiagonalité discrète $T_3$ du crible, qui est elle-même
l'expression la plus simple d'une matrice stochastique non triviale
$2 \times 2$. La balance détaillée discrète et l'auto-adjonction
continue se traduisent l'une l'autre par la transformation
unitaire~\eqref{eq:U_unitaire}.

\subsection{Unification}
\label{sec:localisation_unification}

Les quatre mécanismes ci-dessus manifestent quatre facettes d'une même
dynamique de persistance, avec sa mesure invariante de Haar multiplicative
et sa réversibilité $\Ttwo$ (doublement stochastique). Les liens
M1$\leftrightarrow$M2 (mesure de Liouville $\to$ Haar multiplicative) et
M3$\leftrightarrow$M4 (balance détaillée Gibbs $\to$ antipériodicité
$T_3$) sont explicitement construits ; le lien complet
M1$\leftrightarrow$M2$\leftrightarrow$M3$\leftrightarrow$M4 au sens
d'une équivalence mathématique stricte demande une formulation
catégorielle non incluse ici. Les trois rôles du $\sper = 1/2$ recensés
en \S\ref{sec:parametre} (arithmétique $\Tone$, géométrique $\Tfive$,
spectral de Haar) sont opératoriellement identifiés par cette section.
C'est l'apport structurel principal du programme.

\begin{table}[ht]
\centering
\begin{tabular}{@{}lll@{}}
\toprule
\textbf{Mécanisme} & \textbf{Description} & \textbf{Origine PT}\\
\midrule
M1 — Cellule symplectique & $u_\star\,p_\star = 2\pi$ & $[u, p_u] = i$ (Liouville)\\
M2 — Jacobien de Haar & $du = dp/p$, shift $+1/2$ & $\Ttwo$, doublement stochastique\\
M3 — Transformation unitaire & $U = e^{-su/2}$ & balance détaillée Gibbs $\to$ AA\\
M4 — Condition aux bords & $\theta = \pi$ antipériodique & $T_3$, antidiagonale du crible\\
\bottomrule
\end{tabular}
\caption{Les quatre mécanismes équivalents pour la ligne critique
$\Reop(s) = 1/2$.}
\label{tab:quatre_mecanismes}
\end{table}

